% 第七章 估值分析

\section{估值方法论}

功率半导体行业的估值需要结合行业属性、公司模式和发展阶段综合判断，单一指标往往不够全面。我们建议采用多维度估值框架：

\begin{enumerate}[leftmargin=2em]
\item \textbf{PE 估值}：适用于盈利稳定的成熟型公司，如 IDM 龙头
\item \textbf{PS 估值}：适用于处于投入期、尚未盈利或盈利处于周期底部的公司，如 SiC/GaN 厂商
\item \textbf{PB 估值}：适用于重资产制造型企业，与产能规模高度相关
\item \textbf{EV/EBITDA}：适用于重资产、折旧占比高的公司，反映真实盈利水平
\end{enumerate}

对于功率半导体行业，我们认为 PE 仍是最核心的估值指标，但需注意两个问题：一是行业具有周期性，PE 在周期顶部偏低、周期底部偏高，需要穿越周期看盈利中枢；二是不同业务模式的公司估值差异大，IDM 公司和 Fabless 公司不可简单类比。

\section{历史估值区间}

回顾 2019–2026 年中国功率半导体板块的估值走势，呈现明显的周期性特征：

\begin{itemize}[leftmargin=2em]
\item \textbf{2019–2020 年}：估值低位，行业低谷期 PE 约 30–50 倍
\item \textbf{2020–2021 年}：估值快速扩张，牛市峰值 PE 达到 100–150 倍
\item \textbf{2022 年}：估值回调，PE 回落至 50–70 倍
\item \textbf{2023–2024 年}：行业下行周期，盈利下滑，PE 被动抬升至 80–120 倍（分母端下滑）
\item \textbf{2025–2026 年}：行业复苏，盈利修复，PE 向 50–70 倍回归
\end{itemize}

\begin{exhibitbox}[图 7-1：功率半导体板块历史 PE 走势示意图]
\centering
\vspace{0.3cm}
\begin{tikzpicture}
\begin{axis}[
    width=12cm, height=5.5cm,
    xlabel={年份},
    ylabel={PE (倍)},
    xtick={2019,2020,2021,2022,2023,2024,2025,2026},
    x tick label style={font=\scriptsize},
    ymin=0, ymax=160,
    ymajorgrids=true,
    grid style={dashed, gray!30},
    legend pos=north east,
    legend style={font=\footnotesize},
    no marks,
    thick,
]
% PE走势
\addplot[navy, smooth] coordinates {
    (2019, 40) (2020, 65) (2021, 130) (2022, 55)
    (2023, 90) (2024, 110) (2025, 85) (2026, 70)
};
\legend{板块 PE (TTM)}

% 估值区间带
\draw[dashed, riskred, thick] (2019, 100) -- (2026, 100) node[right, font=\scriptsize, color=riskred] {高估阈值};
\draw[dashed, riskamber, thick] (2019, 60) -- (2026, 60) node[right, font=\scriptsize, color=riskamber] {中枢位置};
\draw[dashed, riskgreen, thick] (2019, 35) -- (2026, 35) node[right, font=\scriptsize, color=riskgreen] {低估阈值};

\end{axis}
\end{tikzpicture}
\vspace{0.2cm}
\sourcenote{Wind，本研究团队整理（示意图，走势经平滑处理）}
\end{exhibitbox}

\section{全球估值对比：中外估值溢价分析}

国内外功率半导体公司的估值差异较大，中国公司普遍存在显著的估值溢价。

\begin{exhibitbox}[表 7-1：国内外功率半导体公司估值对比]
\centering
\small
\begin{tabularx}{\textwidth}{X c c c c}
\toprule
\textbf{公司} & \textbf{市值} & \textbf{PE(TTM)} & \textbf{PS(TTM)} & \textbf{PB} \\
\midrule
\rowcolor{lightgrey}
\multicolumn{5}{l}{\textit{海外龙头}} \\
英飞凌 (Infineon) & ~€450亿 & 18–22x & 3.5–4.0x & 3.0–3.5x \\
意法半导体 (ST) & ~\$400亿 & 16–20x & 2.8–3.2x & 2.5–3.0x \\
安森美 (Onsemi) & ~\$350亿 & 15–18x & 3.8–4.2x & 3.0–3.5x \\
Wolfspeed & ~\$100亿 & 亏损 & 2.5–3.5x & 2.0–2.5x \\
\midrule
\rowcolor{lightgrey}
\multicolumn{5}{l}{\textit{国内龙头}} \\
时代电气 & ~¥804亿 & ~20x & ~2.8x & ~2.3x \\
华润微 & ~¥1,065亿 & ~117x & ~10.5x & ~5.0x \\
士兰微 & ~¥753亿 & ~164x & ~6.7x & ~3.5x \\
斯达半导 & ~¥314亿 & ~95x & ~9.3x & ~7.5x \\
扬杰科技 & ~¥692亿 & ~55x & ~11.5x & ~8.0x \\
三安光电 & ~¥878亿 & 亏损 & ~4.9x & ~1.8x \\
天岳先进 & ~¥724亿 & 亏损 & ~50x & ~6.5x \\
新洁能 & ~¥315亿 & ~83x & ~10.0x & ~8.5x \\
东微半导 & ~¥75亿 & ~35x & ~6.5x & — \\
晶丰明源 & ~¥195亿 & ~247x & ~10.0x & ~7.0x \\
\bottomrule
\end{tabularx}
\sourcenote{Wind, Bloomberg，截至 2026 年 6 月，估值数据因股价波动可能变化}
\end{exhibitbox}

中外估值差异的核心原因：

\begin{enumerate}[leftmargin=2em]
\item \textbf{成长性差异}：中国功率半导体公司正处于国产替代的快速成长期，增速远高于海外成熟龙头，高估值对应高成长预期。
\item \textbf{业务结构差异}：海外龙头以车规级、工业级等高端产品为主，利润率高但增速慢；国内厂商产品结构从中低端向高端升级，利润率较低但增速快。
\item \textbf\textbf{市场环境差异}：A 股市场成长股估值溢价普遍较高，半导体板块享受估值溢价。
\item \textbf{国产替代预期}：市场给予国产替代赛道较高的估值溢价，尤其是 SiC、车规等突破预期强的领域。
\end{enumerate}

\section{估值分层体系}

我们将功率半导体公司按照商业模式、技术壁垒、成长阶段划分为四个估值层级：

\begin{exhibitbox}[表 7-2：功率半导体公司估值分层体系]
\centering
\small
\begin{tabularx}{\textwidth}{l X c L{3.5cm}}
\toprule
\textbf{估值层级} & \textbf{代表性公司} & \textbf{合理PE区间} & \textbf{估值逻辑} \\
\midrule
\rowcolor{navy!10}
\textbf{第一层：成熟IDM龙头} & 时代电气、华润微（复苏后） & 25–40x & 盈利稳定 + 稳健成长 + 制造壁垒 \\
\midrule
\rowcolor{accentblue!10}
\textbf{第二层：设计/IDM成长股} & 斯达半导、新洁能、扬杰 & 40–70x & 高成长 + 国产替代 + 份额提升 \\
\midrule
\rowcolor{riskamber!10}
\textbf{第三层：高景气弹性股} & 东微半导、晶丰明源 & 30–60x & 细分赛道龙头 + 高弹性 \\
\midrule
\rowcolor{steelgrey!20}
\textbf{第四层：前沿技术主题} & 天岳先进、三安光电（化合物半导） & PS 5–20x & 技术突破 + 远期空间 + 尚未盈利 \\
\bottomrule
\end{tabularx}
\sourcenote{本研究团队构建}
\end{exhibitbox}

需要注意的是，估值分层不是静态的，而是随着公司发展阶段变化而动态调整。例如，随着 SiC 业务放量盈利，天岳先进有望从第四层向第三层甚至第二层迁移，带来估值体系的切换。

\section{本研究估值模型精要（AStock House Model）}

为确保估值结论的可复现性与可审计性，本节明确定义本研究采用的估值模型、参数口径与计算流程。所有最终目标价（见第十一章）均由本节参数严格推导得出，不存在未披露的"黑箱"判断。

\subsection{模型定义与优先级}

本研究采用\textbf{相对估值为主、分部估值（SOTP）为辅、风险折价修正}的三层框架。六种估值方法的应用优先级与适用条件如下：

\begin{exhibitbox}[表 7-M1：估值模型——方法优先级与适用范围]
\centering
\small
\begin{tabularx}{\textwidth}{>{\hsize=0.8\hsize}X >{\hsize=0.7\hsize}X >{\hsize=0.9\hsize}X >{\hsize=1.8\hsize}X >{\hsize=1.2\hsize}X}
\toprule
\textbf{方法} & \textbf{优先级} & \textbf{适用对象} & \textbf{核心参数区间} & \textbf{不适用条件} \\
\midrule
\rowcolor{navy!8}
\textbf{PE（2026E）} & 1（核心） & 盈利公司 8/10（除天岳、三安） & 按分层体系（表7-M2）取区间 & 微利/亏损、周期底部净利异常 \\
\midrule
\textbf{PS（2026E）} & 2（亏损替代） & 微利/亏损 2/10（天岳、三安） & Tier4：PS 8–20x & 盈利稳定公司（PS信息含量低） \\
\midrule
\textbf{SOTP} & 3（特殊） & 时代电气（轨交+功率）、三安（LED+SiC） & 各分部独立估值后加总 & 单一业务公司 \\
\midrule
\textbf{PEG（2026E）} & 4（辅助交叉验证） & 成长期公司（斯达、扬杰、新洁能） & 合理 PEG 1.2–1.8x；SiC主题可到 2.0–3.0x & 增速为负或低基数效应 \\
\midrule
\textbf{PB / EV/EBITDA} & 5（周期修正） & 重资产 IDM（华润微、士兰微） & PB 3–6x；EV/EBITDA 15–25x & 轻资产 Fabless \\
\midrule
\textbf{DCF} & 6（敏感性参考） & 不推荐作为主方法 & WACC 8–12\%；终值增长 2–3\% & 周期性+技术迭代型行业（默认不使用） \\
\bottomrule
\end{tabularx}
\sourcenote{本研究团队定义（AStock House Valuation Model v1.0）}
\end{exhibitbox}

关于 DCF 的说明：功率半导体行业兼具强周期性与技术迭代不确定性，DCF 对终值假设与 WACC 的敏感性过高，容易产生"伪精确"。本研究不将 DCF 作为主估值方法，仅在情景分析中作为边际参考；国际投行（Goldman Sachs、Morgan Stanley）对同类周期成长股通常也采用 PE/PEG/SOTP 主框架，与我们一致。

\subsection{分层估值倍数参数表（三情景）}

估值倍数的选取采用「公司分层 × 市场情景」二维矩阵。分层依据公司商业模式、盈利能力、成长确定性（见表 7-2）；情景参数与第十一章三情景假设联动（乐观 25\% / 中性 55\% / 悲观 20\%）。

\begin{exhibitbox}[表 7-M2：估值模型——分层倍数参数（悲观 / 中性 / 乐观）]
\centering
\small
\begin{tabularx}{\textwidth}{l >{\hsize=1.3\hsize}X c c c}
\toprule
\textbf{层级} & \textbf{代表公司} & \textbf{悲观情景 PE/PS} & \textbf{中性情景 PE/PS} & \textbf{乐观情景 PE/PS} \\
\midrule
\rowcolor{navy!8}
\textbf{Tier 1 成熟IDM} & 时代电气（PE）、华润微（复苏后PE） & 15–22x PE / 2.0–2.5x PS & 22–28x PE / 2.6–3.2x PS & 28–38x PE / 3.2–4.0x PS \\
\midrule
\rowcolor{accentblue!8}
\textbf{Tier 2 成长型} & 扬杰科技、斯达半导、新洁能 & 30–40x PE / 6–8x PS & 40–55x PE / 8–12x PS & 55–75x PE / 12–16x PS \\
\midrule
\rowcolor{riskamber!8}
\textbf{Tier 3 高Beta} & 东微半导、士兰微（复苏前高PE用PB修正）、晶丰明源 & 25–35x PE / 4–6x PS & 35–50x PE / 6–9x PS & 50–70x PE / 9–13x PS \\
\midrule
\rowcolor{steelgrey!16}
\textbf{Tier 4 前沿/亏损} & 天岳先进（PS）、三安光电（PB+SOTP） & 6–10x PS / 2.0–2.3x PB & 10–15x PS / 2.3–2.8x PB & 15–22x PS / 2.8–3.5x PB \\
\bottomrule
\end{tabularx}
\sourcenote{本研究团队构建。区间来源：悲观=近3年估值分位 15–25\%；中性=券商一致目标价隐含估值中枢；乐观=2020–21 景气上行期估值上限}
\end{exhibitbox}

\subsection{估值流程（标准五步）}

本研究对每一家公司均执行以下五步，确保估值结论可追溯：

\vspace{0.3em}

\begin{enumerate}[leftmargin=2em]
\item \textbf{Step 1 确定公司分层}（表 7-2 / 表 7-M2），选择主估值方法（表 7-M1 优先级）
\item \textbf{Step 2 锚定 2026E EPS / 营收 中枢}（表 7-M3 一致预期中位数；无一致预期时由本团队基于 IR 指引 + 历史增长推算）
\item \textbf{Step 3 代入三情景倍数}（表 7-M2），得到原始目标价区间
\item \textbf{Step 4 特殊公司执行 SOTP}（时代电气、三安光电；详见表 7-M4 / 7-M5）
\item \textbf{Step 5 风险折价修正}（表 7-12），对高风险公司的\textbf{中性}与\textbf{乐观}目标价上限打折
\end{enumerate}

\vspace{0.3em}

以上五步的每一个参数（EPS、倍数、折价率）均在本节以下各表中完整披露，读者可自行复算验证。

\subsection{Step 2 锚定：2026E 业绩预期全表（10 家公司完整覆盖）}

下表汇总本研究采用的 2026 年业绩预测中枢。数据来源优先级：L1 财报基准 $\rightarrow$ L2 IR 公开指引 $\rightarrow$ L3 券商一致预期中位数 $\rightarrow$ L4 本研究团队推算（明确标注）。

\begin{exhibitbox}[表 7-M3：2026E 业绩预测（10 家公司完整覆盖，含 EPS 与营收）]
\centering
\small
\setlength{\tabcolsep}{3pt}
\begin{tabularx}{\textwidth}{>{\hsize=0.9\hsize}X >{\hsize=0.65\hsize}X c c c c c >{\hsize=1.2\hsize}X}
\toprule
\centering\textbf{公司} & \centering\textbf{总股本} & \textbf{2025A 归母} & \textbf{2026E 归母} & \textbf{2026E} & \textbf{2026E} & \textbf{2026E 营收} & \centering\arraybackslash\textbf{数据来源} \\
\centering (代码) & \centering (亿股) & \textbf{净利(亿)} & \textbf{净利(亿)} & \textbf{增速} & \textbf{EPS(元)} & \textbf{中位数(亿)} & \\
\midrule
\rowcolor{riskgreen!5}
\centering 时代电气 \newline \scriptsize(688187) & \centering ~13.7 & 40.97 & \textbf{46.5} & +13.5\% & ~3.42 & ~310 & 中信建投、华鑫；IR指引（L2）\\
\midrule
\rowcolor{riskgreen!5}
\centering 扬杰科技 \newline \scriptsize(300373) & \centering ~5.42 & 12.59 & \textbf{16.0} & +27.1\% & 2.95 & ~91 & 国金/国信/东海 3家一致（L3）\\
\midrule
\rowcolor{riskamber!5}
\centering 华润微 \newline \scriptsize(688396) & \centering ~14.47 & 6.61 & \textbf{10.85} & +64.1\% & 0.75 & ~128 & 券商一致中位（L3）；IR指引经营利润30\%+\\
\midrule
\rowcolor{riskamber!5}
\centering 士兰微 \newline \scriptsize(600460) & \centering ~16.46 & 3.99 & \textbf{8.40} & +110\% & 0.51 & ~156 & 群益/国信；IR涨价周期确认（L2）\\
\midrule
\rowcolor{riskamber!5}
\centering 斯达半导 \newline \scriptsize(603290) & \centering ~2.40 & 4.05 & \textbf{5.50} & +35.8\% & ~2.29 & ~45 & 财务数据已补全*（L1）；券商一致（L3）\\
\midrule
\rowcolor{riskamber!5}
\centering 新洁能 \newline \scriptsize(605111) & \centering ~4.24 & 3.94 & \textbf{5.00} & +26.9\% & ~1.18 & ~22 & Wind一致推算（L3）\\
\midrule
\rowcolor{riskred!5}
\centering 天岳先进 \newline \scriptsize(688234) & \centering ~4.85 & -2.08 & \textbf{0.51（扭亏）} & 扭亏为盈 & 0.11 & ~20.6 & 华源/国信；SiC价格触底确认（L3）\\
\midrule
\rowcolor{steelgrey!10}
\centering 三安光电 \newline \scriptsize(600703) & \centering ~49.9 & -3.53 & \textbf{~0（减亏）} & 大幅减亏 & ~0 & ~195–200 & 困境反转区间估计（L4\textsuperscript{\textdagger}）\\
\midrule
\rowcolor{riskamber!5}
\centering 东微半导 \newline \scriptsize(688261) & \centering ~1.01 & 0.46 & \textbf{0.85} & +84.8\% & ~0.84 & ~15.5 & 本团队推算（L4\textsuperscript{\textdagger}）：高压MOS修复+Q1低基数\\
\midrule
\rowcolor{riskamber!5}
\centering 晶丰明源 \newline \scriptsize(688368) & \centering ~0.89 & 0.36 & \textbf{1.75} & +386\% & ~1.97 & ~22 & 本团队推算（L4\textsuperscript{\textdagger}）：AI PMIC爆发，盈利预测近3月上调80\%\\
\bottomrule
\end{tabularx}
\vspace{0.2cm}
\sourcenote{L1=公司定期报告原文；L2=IR公开交流指引；L3=券商研报一致预期中位数；L4=本团队推算\\
*斯达半导 2025A 归母净利润 4.05 亿来自定期报告（official\_financials L32），此前审计表 7-9 因数据版本滞后标注"缺失"，现已补全验证；修正后 PE(TTM)=77.5x（TTM口径）\\
\textsuperscript{\textdagger} 华润微 IR 指引口径为"经营利润进一步增长 30\%+"，与归母净利增速 64\% 差异源于非经常性损益与12吋折旧递减结构，详见下方口径澄清\\
\textsuperscript{\textdagger} 东微半导与晶丰明源缺乏公开券商覆盖，为本团队基于26Q1年化+产品结构+订单趋势推算，较其他标的预测不确定性高}
\end{exhibitbox}

\textbf{关于华润微增速的口径澄清}：华润微管理层 IR 指引"经营利润 30\%+ 增长"与本研究采用的归母净利 +64.1\%（券商一致 10.85 亿）并不矛盾。差异来自三方面：（1）2025 年归母净利含股权激励与资产减值等非经常性损失约 1.2 亿，2026 年基数回归正常；（2）重庆 12 吋产线折旧从 2025 年高峰期进入递减期，对归母净利释放正贡献；（3）"经营利润"为扣非口径，通常低于归母净利增速。取券商一致 10.85 亿为中枢具备市场可验证性。

\subsection{Step 4 特殊公司：SOTP 分部估值（时代电气 + 三安光电）}

两家公司业务高度多元化，单一 PE/PS 会产生系统性估值偏置（详见估值审计 §3.3）。本研究对其执行 SOTP。

\begin{exhibitbox}[表 7-M4：时代电气（688187）SOTP 分部估值]
\centering
\small
\begin{tabularx}{\textwidth}{X c c c c c}
\toprule
\textbf{业务分部} & \textbf{2026E 营收(亿)} & \textbf{占比} & \textbf{2026E 净利(亿)} & \textbf{分部估值方法} & \textbf{分部估值(亿)} \\
\midrule
\rowcolor{navy!5}
\textbf{轨道交通装备（含新产业）} & ~195 & ~63\% & ~29.5 & PE 18–22x（A股轨交装备可比：中国中车16x、交控科技22x） & 531–649 \\
\midrule
\rowcolor{accentblue!5}
\textbf{功率半导体（含 SiC）} & ~85 & ~27\% & ~12.5 & PE 30–40x（Tier 1 IDM：华润微可比上沿） & 375–500 \\
\midrule
\rowcolor{riskamber!5}
\textbf{新能源电驱/储能/传感器} & ~30 & ~10\% & ~4.5 & PE 25–35x（汽车零部件可比：汇川技术28x） & 113–158 \\
\midrule
\textbf{合计（加总）} & \textbf{310} & 100\% & \textbf{46.5} & — & \textbf{1,019–1,307} \\
\midrule
\midrule
\multicolumn{5}{r}{\textbf{对应 A 股目标价（总股本 13.7 亿股）}} & \textbf{74.4 – 95.4 元} \\
\multicolumn{5}{r}{\textbf{叠加 H 股折价因子（A/H 溢价 46\%，当前已在价格中反映）}} & 不做额外调整 \\
\bottomrule
\end{tabularx}
\sourcenote{分部利润拆分参考公司业务结构历史占比与IR披露；分部PE选取A股同行业可比公司2026E中位}
\end{exhibitbox}

SOTP 中性估值区间 74–95 元，与第十一章基于整体 Tier 1 PE 法得到的中性目标 75–85 元高度吻合，\textbf{交叉验证通过}。SOTP 的意义在于揭示了时代电气"低估"的结构性来源——功率半导体分部给予 30–40x 后估值贡献已达 375–500 亿，市场当前仅给予整体 804 亿市值，意味着轨交装备分部实际隐含估值仅 ~300 亿（对应 PE ~10x，显著低于同业 16–22x），\textbf{这是时代电气估值折价的数学锚点，而非主观判断}。

\begin{exhibitbox}[表 7-M5：三安光电（600703）SOTP 分部估值]
\centering
\small
\begin{tabularx}{\textwidth}{X c c c c c}
\toprule
\textbf{业务分部} & \textbf{2026E 营收(亿)} & \textbf{占比} & \textbf{盈利状态} & \textbf{分部估值方法} & \textbf{分部估值(亿)} \\
\midrule
\rowcolor{steelgrey!10}
\textbf{LED 芯片（含 Mini/Micro LED）} & ~135 & ~69\% & 微利~盈亏平衡 & PB 2.0–2.5x；净资产对应约 250 亿分部账面值 & 500–625 \\
\midrule
\rowcolor{accentblue!5}
\textbf{化合物半导体（SiC/GaN/射频）} & ~40 & ~20\% & 亏损收窄 & PS 4–7x（Tier 4 下限：纯SiC天岳先进 PS 10–15x 为上界，三安SiC进度落后一档取折价） & 160–280 \\
\midrule
\rowcolor{riskamber!5}
\textbf{光通讯/滤波器/其他} & ~22 & ~11\% & 微利 & PB 2.0x 估算 & 100–140 \\
\midrule
\textbf{合计（加总）} & \textbf{197} & 100\% & 整体减亏 & — & \textbf{760–1,045} \\
\midrule
\midrule
\multicolumn{5}{r}{\textbf{对应目标价（总股本 49.9 亿股）}} & \textbf{15.2 – 20.9 元} \\
\bottomrule
\end{tabularx}
\sourcenote{分部结构参考公司公告与机构调研；化合物半导体参考Tier 4估值体系并打工艺折价}
\end{exhibitbox}

三安光电 SOTP 估值区间 15–21 元与当前股价 ~17.5 元基本持平，反映市场对 LED 业务复苏与 SiC 业务贡献仍持观望态度——\textbf{这也是我们给予"中性"评级的定量依据，而非主观判断"困境反转太慢"}。

\subsection{估值模型的可审计性声明}

本研究所有目标价均严格按以下公式推导，所有参数已在本节与表 7-12 中完整披露，读者可使用电子表格复算：

\begin{align*}
\text{悲观目标价下限} &= \text{2026E EPS} \times \text{该层级悲观倍数下限} \times (1 - \text{风险折价率下限}) \\
\text{中性目标价区间} &= \text{2026E EPS} \times \text{该层级中性倍数区间} \times (1 - \text{风险折价率中位}) \\
\text{乐观目标价上限} &= \text{2026E EPS} \times \text{该层级乐观倍数上限} \times (1 - 0) \quad \text{（乐观情景不叠加折价）}
\end{align*}

对 Tier 4 亏损公司，公式中 EPS 替换为 每股营收（2026E 营收 $\div$ 总股本），PE 替换为 PS。对 SOTP 公司（时代电气、三安），以 SOTP 计算结果为主，整体 PE 法结果作为交叉验证边界。

\vspace{0.3em}

\section{估值水平判断}

基于以上分析，我们对当前功率半导体板块的估值水平判断如下：

\begin{keyinsight}[估值水平整体判断]
当前功率半导体板块整体估值处于历史中位偏高的水平，但分化严重——传统成熟型公司（如时代电气）估值处于历史低位，高成长赛道公司（如 SiC、AI 相关）估值处于历史中高位。

投资建议上，优先选择估值合理、业绩确定性高的龙头公司（如时代电气、扬杰科技），同时战略性配置第三代半导体等高成长赛道。对于估值过高、业绩兑现不确定性大的公司，建议保持谨慎。
\end{keyinsight}

从纵向（历史对比）和横向（全球对比）两个维度看：

\begin{itemize}[leftmargin=2em]
\item \textbf{纵向看}：板块整体 PE 约 80–100 倍，处于历史中位偏高区间，但需注意当前盈利处于周期底部，随着业绩复苏 PE 将被动回落。
\item \textbf{横向看}：国内公司估值显著高于海外龙头，但考虑到成长性差异（国内增速 2–3 倍于海外），估值溢价有一定合理性。
\end{itemize}

我们认为，判断估值是否合理的关键在于\textbf{业绩增长的可持续性和确定性}。如果国产替代和 AI/新能源双轮驱动的逻辑能够持续兑现，当前估值仍有支撑；反之，如果需求不及预期或竞争加剧，估值存在回调风险。
% 第七章估值分析扩展：机构目标价、业绩兑现与估值审计
% 本文件为 ch07_valuation.tex 的补充章节，需在 \section{估值水平判断} 之后引入

\section{机构一致预期与目标价历史}

卖方机构的盈利预测与目标价是市场预期的集中体现，也是估值分析的重要参照。我们梳理了覆盖八家重点功率半导体公司的主要券商研报，从评级分布、目标价区间、盈利预测分歧三个维度进行分析。需要强调的是，卖方目标价存在系统性向上偏差（通常高估 20\%–40\%），实际投资决策中建议打 7–8 折作为参考。

\begin{exhibitbox}[表 7-3：卖方评级分布与覆盖度统计]
\centering
\small
\begin{tabularx}{\textwidth}{X c c c c c c}
\toprule
\textbf{公司} & \textbf{买入} & \textbf{增持/优于大市} & \textbf{中性} & \textbf{一致评级} & \textbf{覆盖机构数} & \textbf{近3月上调次数} \\
\midrule
时代电气 & 8 & 2 & 0 & 买入 & 10 & 1 \\
扬杰科技 & 6 & 3 & 1 & 买入 & 10 & 2 \\
斯达半导 & 5 & 4 & 1 & 买入/增持 & 10 & 0 \\
新洁能 & 3 & 5 & 2 & 增持 & 10 & 0 \\
华润微 & 3 & 5 & 2 & 增持 & 10 & 1 \\
士兰微 & 2 & 6 & 2 & 增持 & 10 & 2 \\
天岳先进 & 2 & 3 & 5 & 中性/增持 & 10 & 0 \\
三安光电 & 1 & 4 & 5 & 中性 & 10 & 0 \\
\bottomrule
\end{tabularx}
\sourcenote{Wind 一致预期、各券商研报统计（近 90 天，截至 2026 年 6 月）}
\end{exhibitbox}

从评级分布看，板块内部分化明显。时代电气与扬杰科技获得最多"买入"评级，反映机构对其业绩确定性的认可；而天岳先进与三安光电的评级分歧较大，中性评级占比过半，体现市场对其盈利前景的疑虑。近三个月的评级上调集中在行业涨价受益的 IDM 龙头（士兰微、扬杰科技、华润微），显示机构对行业复苏方向的共识正在形成。

\begin{exhibitbox}[表 7-4：一致目标价区间与隐含涨跌空间]
\centering
\small
\begin{tabularx}{\textwidth}{X c c c c c}
\toprule
\textbf{公司} & \textbf{当前股价(元)} & \textbf{一致目标价区间(元)} & \textbf{目标价区间(元)} & \textbf{隐含涨幅} & \textbf{分歧度*} \\
\midrule
\rowcolor{riskgreen!10}
时代电气 & ~58.7 & 75–85 & 68–95 & +27\%~+44\% & 低 \\
\rowcolor{riskgreen!10}
士兰微 & ~41.8 & 50–55 & 45–65 & +20\%~+31\% & 中低 \\
\rowcolor{riskamber!10}
扬杰科技 & ~123.1 & 150–170 & 130–200 & +22\%~+38\% & 中低 \\
\rowcolor{riskamber!10}
斯达半导 & ~130.9 & 150–180 & 120–220 & +15\%~+37\% & 中 \\
\rowcolor{riskamber!10}
华润微 & ~73.6 & 85–100 & 70–120 & +15\%~+35\% & 中 \\
\rowcolor{riskamber!10}
新洁能 & ~74.3 & 90–110 & 75–130 & +22\%~+49\% & 中高 \\
\rowcolor{riskred!10}
天岳先进 & ~149.2 & 180–220 & 120–300 & +21\%~+47\% & \textbf{高} \\
\rowcolor{riskred!10}
三安光电 & ~17.5 & 20–25 & 15–30 & +14\%~+43\% & \textbf{高} \\
\bottomrule
\end{tabularx}
\sourcenote{各券商目标价汇总，Wind 一致预期（截至 2026 年 6 月 21 日）\\
*分歧度 = (最高目标价 - 最低目标价) / 一致目标价中位数；低：<20\%，中低：20–30\%，中：30–50\%，中高：50–70\%，高：>70\%}
\end{exhibitbox}

目标价的分歧度与公司的业绩确定性高度相关。时代电气分歧度最低，因为轨交基本盘稳健，盈利可预测性强；天岳先进和三安光电分歧度最高，反映市场对 SiC 周期拐点和 LED 复苏节奏的巨大认知差异。值得注意的是，若按卖方目标价打 7 折后的"实际可实现目标价"计算，天岳和三安的上涨空间将显著收窄，甚至可能与当前股价持平，这是高分歧标的需要警惕的风险。

接下来我们考察盈利预测的一致性。

\begin{exhibitbox}[表 7-5：2026E 盈利预测对比与机构分歧]
\centering
\small
\begin{tabularx}{\textwidth}{X c c c c c c}
\toprule
\textbf{公司} & \textbf{最高(亿元)} & \textbf{最低(亿元)} & \textbf{中位数(亿元)} & \textbf{最高/最低} & \textbf{调整方向(近3月)} & \textbf{数据来源数} \\
\midrule
\rowcolor{riskgreen!10}
时代电气 & ~47.0 & ~45.8 & ~46.5 & 1.03x & 上调(小幅) & 2+ \\
\rowcolor{riskgreen!10}
扬杰科技 & 16.68 & 15.91 & 16.0 & 1.05x & 上调 & 3 \\
\rowcolor{riskamber!10}
士兰微 & 9.48 & 8.09 & 8.4 & 1.17x & 上调 & 3 \\
\rowcolor{riskamber!10}
新洁能 & ~5.5 & ~4.5 & ~5.0 & 1.22x & 上调 & 1+ \\
\rowcolor{riskamber!10}
华润微 & 12.00 & 9.70 & 10.85 & 1.24x & 上调(小幅) & 2 \\
\rowcolor{riskamber!10}
斯达半导 & ~6.0 & ~4.5 & ~5.5 & 1.33x & 不确定 & 1+ \\
\rowcolor{riskred!10}
天岳先进 & 0.63 & 0.38 & 0.51 & 1.66x & \textbf{下调(大幅)} & 2 \\
\rowcolor{steelgrey!10}
三安光电 & 微利 & 亏损 & ~0 & — & 减亏 & 1+ \\
\bottomrule
\end{tabularx}
\sourcenote{国信证券、华源证券、华鑫证券、东海证券、国金证券、群益证券等公开研报及 Wind 一致预期}
\end{exhibitbox}

盈利预测的分歧度验证了前述判断：时代电气和扬杰科技的预测高度一致（最高/最低比低于 1.1x），而天岳先进的分歧超过 1.6x。近三个月的预测调整方向呈现明显的"IDM 上调、设计/第三代分化"格局——华润微、士兰微、扬杰科技等 IDM 厂商因 AI 需求驱动和行业涨价周期获得盈利预测上调，而斯达半导、天岳先进等或因价格战、或因周期位置不明朗面临下调压力。

\begin{keyinsight}[机构预期的投资含义]
综合评级分布、目标价分歧度和盈利预测调整方向，我们得出以下投资启示：

1. \textbf{低分歧+上调趋势}的公司（时代电气、扬杰科技）盈利确定性最高，适合作为核心配置，预期收益虽不夸张但胜率高。
2. \textbf{高分歧但方向正在上修}的公司（士兰微、华润微）存在预期差修复机会，适合弹性配置，但需警惕周期波动。
3. \textbf{高分歧+下调趋势}的公司（天岳先进、斯达半导）左侧抄底风险大，建议等待业绩拐点明确后再介入。
4. 卖方目标价整体存在 20\%–40\% 的系统性高估，实际操作中建议打 7–8 折作为合理估值上限。
\end{keyinsight}

\section{业绩预期与兑现追踪}

估值分析不能脱离业绩兑现能力。我们追踪了九家重点公司近四个季度的业绩表现，对比市场一致预期，构建 Beat/Miss 矩阵，以评估各家公司的业绩兑现质量和预期管理能力。业绩兑现的持续性是支撑估值的核心因素——持续超预期的公司往往享受估值溢价，而频繁低于预期的公司则面临估值下杀。

\subsection{近四季度 Beat/Miss 矩阵}

我们以归母净利润偏离一致预期 ±5\% 为界，将业绩表现分为 Beat（超预期）、In-line（符合预期）、Miss（不及预期）三档。对于亏损公司，则以营收和毛利率为主要判断依据。

\begin{exhibitbox}[表 7-6：近四季度业绩 Beat/Miss 矩阵]
\centering
\small
\begin{tabularx}{\textwidth}{>{\hsize=0.6\hsize}X >{\hsize=1.0\hsize}X c >{\hsize=1.0\hsize}X c >{\hsize=1.4\hsize}X}
\toprule
\centering\textbf{公司} & \centering\textbf{2025全年} & \textbf{2025Q4} & \centering\textbf{2026Q1} & \textbf{趋势} & \centering\arraybackslash\textbf{近4季综合判断} \\
\midrule
\rowcolor{riskgreen!10}
\centering 晶丰明源 & \centering Beat（扭亏超预期） & Beat & \centering Beat（大幅） & \textcolor{riskgreen}{\textbf{持续↑}} & \centering\arraybackslash 持续 Beat \\
\rowcolor{riskgreen!10}
\centering 扬杰科技 & \centering In-line（稳健增长） & Miss（Q4单季） & \centering Beat（小幅） & \textcolor{riskgreen}{\textbf{改善↑}} & \centering\arraybackslash 偏 Beat \\
\rowcolor{riskamber!10}
\centering 华润微 & \centering In-line（周期底部） & In-line & \centering Beat（大幅） & \textcolor{riskgreen}{\textbf{改善↑}} & \centering\arraybackslash 明显改善，转 Beat \\
\rowcolor{riskamber!10}
\centering 士兰微 & \centering In-line（减亏改善） & — & \centering In-line~Beat & \textcolor{riskgreen}{\textbf{改善↑}} & \centering\arraybackslash 温和改善 \\
\rowcolor{riskamber!10}
\centering 天岳先进 & \centering Miss（大幅，价格战） & Miss & \centering Beat（毛利率改善） & \textcolor{riskgreen}{\textbf{反转↑}} & \centering\arraybackslash 底部反转，从 Miss 转 Beat \\
\rowcolor{riskamber!10}
\centering 时代电气 & \centering In-line（稳健） & In-line & \centering Miss（半导体板块） & \textcolor{riskred}{\textbf{承压↓}} & \centering\arraybackslash 分化，整体 In-line \\
\rowcolor{riskred!10}
\centering 新洁能 & \centering In-line（增速放缓） & — & \centering In-line~Miss & \textcolor{riskred}{\textbf{偏弱→}} & \centering\arraybackslash 偏弱，接近 Miss \\
\rowcolor{riskred!10}
\centering 斯达半导 & \centering Miss（增收不增利） & — & \centering Miss（显著） & \textcolor{riskred}{\textbf{持续↓}} & \centering\arraybackslash 持续 Miss \\
\rowcolor{riskred!10}
\centering 三安光电 & \centering Miss（大幅，由盈转亏） & Miss & \centering Miss（显著） & \textcolor{riskred}{\textbf{持续↓}} & \centering\arraybackslash 持续 Miss \\
\bottomrule
\end{tabularx}
\sourcenote{公司定期报告、各券商财报点评、Wind 一致预期（数据区间：2025Q2–2026Q1）}
\end{exhibitbox}

Beat/Miss 矩阵清晰地揭示了板块内部的业绩分化格局。\textbf{持续 Beat}型公司仅有晶丰明源一家，AI 电源芯片的爆发式增长持续超出市场预期；\textbf{改善型}公司（扬杰科技、华润微、士兰微）占据主流，反映行业整体处于复苏通道；\textbf{持续 Miss}型公司（斯达半导、三安光电）则面临各自的结构性问题——IGBT 价格战和 LED 业务拖累。

\subsection{业绩质量与兑现能力分析}

业绩兑现不仅要看"有没有达到预期"，更要看"用什么方式达到预期"。我们从收入利润匹配度、毛利率趋势、现金流质量三个维度评估业绩质量。

\begin{exhibitbox}[表 7-7：业绩质量多维度对比]
\centering
\small
\begin{tabularx}{\textwidth}{X c c c c}
\toprule
\textbf{公司} & \textbf{营收利润匹配度} & \textbf{毛利率趋势} & \textbf{现金流质量} & \textbf{综合质量评价} \\
\midrule
\rowcolor{riskgreen!10}
扬杰科技 & 利润>收入，良性 & 持续上升（36.8\%） & 优秀（净现比 1.35） & \textbf{高质量} \\
\rowcolor{riskgreen!10}
华润微 & 低基数下利润弹性大 & 底部企稳（25.5\%） & 优秀（净现比 3.15） & \textbf{改善中} \\
\rowcolor{riskamber!10}
士兰微 & 利润>收入 & 触底回升（19.8\%） & 良好（净现比 3.75） & \textbf{改善中} \\
\rowcolor{riskamber!10}
晶丰明源 & 扭亏阶段，利润弹性大 & 高位稳定（37.4\%） & 波动大（净现比 5.0） & \textbf{改善中} \\
\rowcolor{riskamber!10}
时代电气 & 收入>利润，略有背离 & 基本稳定（33.3\%） & 良好（净现比 0.97） & \textbf{有压力} \\
\rowcolor{riskred!10}
新洁能 & 收入微增利润降 & 持续下滑（28.8\%） & 良好（净现比 0.98） & \textbf{偏弱} \\
\rowcolor{riskred!10}
斯达半导 & 严重背离（双降） & 持续大幅下滑（21.1\%） & 一般 & \textbf{差} \\
\rowcolor{riskred!10}
三安光电 & 严重背离 & 低位改善（18.5\%） & 良好 & \textbf{差} \\
\rowcolor{riskred!10}
天岳先进 & 双降（2025），边际改善 & V型反转（19.1\%） & 一般 & \textbf{底部} \\
\bottomrule
\end{tabularx}
\sourcenote{公司定期报告，本研究团队整理（毛利率为 2026Q1 数据）}
\end{exhibitbox}

从业绩质量看，\textbf{扬杰科技}是唯一实现"营收利润双增+毛利率持续上升+现金流优秀"的高质量增长公司，这与其 IDM 模式+产品结构升级的战略直接相关。华润微和士兰微虽然当前毛利率仍处于低位，但边际改善趋势明确，且 IDM 模式下的现金流质量显著优于账面利润，属于典型的周期复苏型标的。

值得警惕的是斯达半导和三安光电，营收和利润双降、毛利率持续下滑，业绩质量处于恶化通道，即使估值看起来"便宜"也可能是价值陷阱。

\subsection{一致预期修正与业绩指引}

近三个月的盈利预测调整是市场对业绩兑现结果的直接反馈。我们观察到一个重要的模式：\textbf{业绩超预期的公司，预测上调幅度往往大于超预期幅度本身}——即市场倾向于"外推"好业绩，这会带来估值和盈利的"戴维斯双击"；反之，业绩低于预期的公司，预测下调也会超调。

\begin{exhibitbox}[表 7-8：近 3 个月 2026E 盈利预测调整汇总]
\centering
\small
\begin{tabularx}{\textwidth}{X c c c c}
\toprule
\textbf{公司} & \textbf{3个月前预测(亿元)} & \textbf{当前预测(亿元)} & \textbf{变化幅度} & \textbf{主要驱动因素} \\
\midrule
\rowcolor{riskgreen!10}
晶丰明源 & ~0.8–1.0 & ~1.5–2.0 & \textbf{+80\%}（大幅上调） & AI 电源芯片需求爆发 \\
\rowcolor{riskgreen!10}
天岳先进 & 亏损~1.0 & 0.38（扭亏） & \textbf{大幅上调} & 毛利率大幅改善、价格触底 \\
\rowcolor{riskgreen!10}
华润微 & ~9.7 & ~12.0 & +24\% & Q1 超预期、AI 需求确定 \\
\rowcolor{riskgreen!10}
士兰微 & ~8.0–8.4 & ~9.5 & +15\% & 涨价周期确认、业绩弹性显现 \\
\rowcolor{riskgreen!10}
扬杰科技 & ~14.5 & ~15.9–16.0 & +10\% & 汽车电子超预期、SiC 放量 \\
\rowcolor{riskamber!10}
时代电气 & ~43 & ~45 & +5\%（基本持平） & 稳健增长，预期变化小 \\
\rowcolor{riskred!10}
新洁能 & ~4.5 & ~4.0–4.2 & -10\% & 毛利率下滑、价格竞争 \\
\rowcolor{riskred!10}
斯达半导 & ~5.5–6.0 & ~4.5–5.0 & -15\% & Q1 低于预期、价格战持续 \\
\rowcolor{riskred!10}
三安光电 & 盈利~2.0 & 盈亏平衡~微利 & 大幅下调 & LED 拖累超预期 \\
\bottomrule
\end{tabularx}
\sourcenote{各券商研报盈利预测对比（2026 年 3 月 vs 2026 年 6 月）}
\end{exhibitbox}

预期调整的方向与幅度验证了我们的业绩质量判断——晶丰明源和天岳先进的盈利预测上调幅度最大（均超过 50\%），分别代表"AI 成长"和"周期反转"两种超预期模式；而三安光电和斯达半导的预测持续下调，反映基本面压力尚未释放完毕。

\begin{keyinsight}[业绩兑现将成为估值分化的核心驱动力]
当前板块估值分化已经较为显著，但我们认为\textbf{业绩兑现能力的差异将进一步拉大估值差距}：
1. 持续 Beat + 高质量增长的公司（如扬杰科技），将享受"业绩+估值"双升的戴维斯双击，估值中枢有望上移。
2. 周期反转型公司（如华润微、天岳先进），一旦业绩连续超预期形成趋势，估值修复空间最大，但波动也最大。
3. 持续 Miss 的公司（如斯达半导、三安光电），即使当前估值看似不贵，仍可能面临"业绩下修+估值压缩"的双杀，建议规避或轻仓试探。
4. 稳健型公司（如时代电气），业绩波动小，估值相对稳定，适合作为板块底仓配置。
\end{keyinsight}

\section{估值审计与业绩匹配度分析}

前述估值分析建立在多重假设之上，其可靠性需要通过独立的审计视角进行验证。本节从数学校验、估值方法合理性、可比公司选择、业绩--估值匹配度四个维度对估值体系进行系统性审计，并提出风险调整后的修正估值框架。

\subsection{估值数学校验}

我们对各公司的核心估值指标进行了独立复算，验证数据准确性。

\begin{exhibitbox}[表 7-9：核心估值指标数学校验结果]
\centering
\small
\begin{tabularx}{\textwidth}{X c c c c}
\toprule
\textbf{公司} & \textbf{校验指标} & \textbf{标注值} & \textbf{复算值} & \textbf{差异及说明} \\
\midrule
\rowcolor{riskgreen!10}
时代电气 & PE(TTM) & ~19.6x & 19.6x & 0.0\%，准确 \\
\rowcolor{riskgreen!10}
扬杰科技 & PE(TTM) & ~55x & 51.2x & -7.4\%，TTM 计算时点差异 \\
\rowcolor{riskamber!10}
华润微 & PE(TTM) & ~117x & 117.4x & +0.3\%，准确（TTM 口径） \\
\rowcolor{riskamber!10}
华润微 & 2026E PE & ~74x & 81.9x & +10.7\%，隐含更乐观的盈利预测 \\
\rowcolor{riskamber!10}
士兰微 & PE(TTM) & ~164x & 188.2x & +14.8\%，2025A 口径差异 \\
\rowcolor{riskamber!10}
斯达半导 & PE(TTM) & ~95.7x & ~77.5x & -19.0\%，2025A 归母 4.05 亿已补全*（L1），TTM 计算口径差异 \\
\rowcolor{riskgreen!10}
新洁能 & PE(TTM) & ~82.8x & 83.9x & +1.3\%，接近 \\
\rowcolor{riskgreen!10}
全样本 & PS(TTM) & — & — & 全部 < ±0.5\%，准确 \\
\bottomrule
\end{tabularx}
\sourcenote{本研究团队根据公司财报、市值数据独立复算}
\end{exhibitbox}

总体来看，PS 指标全部准确，PE 指标大部分在可接受误差范围内（±10\% 以内）。需要注意两个问题：一是\textbf{华润微 2026E PE 标注值低于复算值约 11\%}，可能隐含了比公开一致预期更乐观的盈利预测；二是\textbf{斯达半导 PE(TTM) 标注值 95.7x 与按 L1 财报复算的 77.5x 存在约 19\% 差异}，差异来源于 TTM 滚动窗口不同（标注可能使用 Wind 自带 TTM 口径），但 2025A 归母净利 4.05 亿的 L1 数据已从财报确认无误（见 official\_financials L32），原"数据缺失"标注现已消除。

\subsection{估值方法合理性审计}

不同估值方法适用于不同类型的公司，方法选择不当会导致估值结论失真。

\begin{exhibitbox}[表 7-10：各公司估值方法适用性评估]
\centering
\small
\begin{tabularx}{\textwidth}{X c c c c c}
\toprule
\textbf{公司} & \textbf{PE 适用性} & \textbf{PS 适用性} & \textbf{PB 适用性} & \textbf{PEG 适用性} & \textbf{推荐估值方法} \\
\midrule
\rowcolor{riskgreen!10}
时代电气 & 高（盈利稳定） & 中 & 高（重资产） & 中 & \textbf{PE + PB} \\
\rowcolor{riskgreen!10}
扬杰科技 & 高（成长确定） & 中 & 中 & 高（成长期） & \textbf{PE + PEG} \\
\rowcolor{riskamber!10}
华润微 & 中（周期底部） & 中 & 高（IDM） & 低（周期干扰） & \textbf{PB + EV/EBITDA} \\
\rowcolor{riskamber!10}
士兰微 & 低（周期底部PE高） & 中 & 高（IDM） & 中 & \textbf{PB + EV/EBITDA} \\
\rowcolor{riskamber!10}
斯达半导 & 中（盈利下滑） & 中 & 中 & 低（增速为负） & \textbf{PS + PB} \\
\rowcolor{riskred!10}
天岳先进 & 极低（微利/亏损） & 高 & 中 & 极低 & \textbf{PS + SOTP} \\
\rowcolor{riskred!10}
三安光电 & 极低（亏损） & 高 & 高（重资产） & 极低 & \textbf{PB + SOTP} \\
\rowcolor{riskamber!10}
晶丰明源 & 中（高增长） & 中 & 低 & 高（成长期） & \textbf{PE + PEG} \\
\bottomrule
\end{tabularx}
\sourcenote{本研究团队评估}
\end{exhibitbox}

我们的核心发现是：\textbf{当前板块处于周期底部，PE 估值法的参考性显著下降}。对于周期底部的 IDM 公司（华润微、士兰微），PB 和 EV/EBITDA 比 PE 更能反映真实估值水平，因为 PE 在周期底部会因分母变小而被动抬升，造成"贵"的假象。反过来，周期顶部 PE 低但实际估值贵——这是周期股估值的经典陷阱。

\subsection{业绩--估值匹配矩阵}

估值是否合理，最终要回答的问题是：\textbf{业绩增长能否支撑当前估值}？我们构建业绩--估值匹配矩阵，从"业绩兑现质量"和"估值水平"两个维度对公司进行定位。

\begin{exhibitbox}[表 7-11：业绩--估值匹配矩阵]
\centering
\small
\begin{tabularx}{\textwidth}{X c c c c}
\toprule
\textbf{公司} & \textbf{2026E PE} & \textbf{2026E 增速} & \textbf{PEG(2026E)} & \textbf{估值--业绩匹配判断} \\
\midrule
\rowcolor{riskgreen!10}
时代电气 & ~17x & 15–20\% & 0.97–1.13 & \textbf{低估}（PEG < 1） \\
\rowcolor{riskgreen!10}
扬杰科技 & ~42x & 25–30\% & 1.40–1.68 & \textbf{合理}（在合理区间内） \\
\rowcolor{riskamber!10}
士兰微 & ~82x & 130\%+（低基数） & ~0.6 & 看似便宜，但低基数下 PEG 失效 \\
\rowcolor{riskamber!10}
华润微 & ~82x & 60–70\%（低基数） & ~1.1–1.4 & 周期底部，需结合 PB 判断 \\
\rowcolor{riskred!10}
新洁能 & ~55x & 20–30\% & 1.8–2.8 & \textbf{偏高}（PEG 超合理区间） \\
\rowcolor{riskred!10}
斯达半导 & ~57x & 10–20\%（困境反转） & 2.9–5.7 & \textbf{偏高}（反转确定性不足） \\
\rowcolor{riskred!10}
天岳先进 & ~1356x & 扭亏 & 失效 & \textbf{主题投资}，PS ~50x 远高于海外可比 \\
\rowcolor{riskred!10}
三安光电 & 亏损 & 扭亏 & 失效 & \textbf{困境反转}，PB ~2.5x 不算贵但不盈利 \\
\rowcolor{riskamber!10}
晶丰明源 & — & ~100\%+（低基数） & — & 高增长支撑高估值，但可持续性存疑 \\
\bottomrule
\end{tabularx}
\sourcenote{本研究团队测算，增速为 2026E 归母净利润同比增速，PEG 合理区间取 1.2–1.8}
\end{exhibitbox}

按 PEG 框架衡量，板块中\textbf{仅时代电气明确低估，扬杰科技处于合理区间}，其余公司普遍偏贵。但需要注意两个限定条件：一是周期底部 PE 高并不代表估值贵，对于华润微和士兰微这类周期复苏型公司，用正常化盈利计算的 PE 会显著降低；二是天岳先进、三安光电等亏损公司 PE 估值失效，需要用 PS、PB 和 SOTP 等方法评估。

\subsection{风险调整后的估值修正}

前述估值均未考虑风险折价。我们识别出技术路线风险、周期波动风险、地缘政治风险三大核心风险因子，并据此对估值进行风险调整。

\begin{exhibitbox}[表 7-12：风险因子评估与估值折价修正]
\centering
\small
\begin{tabularx}{\textwidth}{X c c c c c}
\toprule
\textbf{公司} & \textbf{技术路线风险} & \textbf{周期波动风险} & \textbf{地缘竞争风险} & \textbf{综合折价率} & \textbf{风险调整后估值} \\
\midrule
\rowcolor{riskgreen!10}
时代电气 & 低 & 低 & 低 & 5–10\% & 基本不变 \\
\rowcolor{riskgreen!10}
扬杰科技 & 低 & 中 & 低 & 10–15\% & 下调 10–15\% \\
\rowcolor{riskamber!10}
华润微 & 低 & 高 & 中 & 20–25\% & 下调 20–25\% \\
\rowcolor{riskamber!10}
士兰微 & 中 & 高 & 中 & 25–35\% & 下调 25–35\% \\
\rowcolor{riskamber!10}
斯达半导 & 中 & 高 & 中 & 20–30\% & 下调 20–30\% \\
\rowcolor{riskred!10}
天岳先进 & 高 & 高 & 高 & 40–50\% & 下调 40–50\% \\
\rowcolor{riskred!10}
三安光电 & 中高 & 中 & 中 & 25–35\% & 下调 25–35\% \\
\rowcolor{riskamber!10}
新洁能 & 低 & 高 & 低 & 15–25\% & 下调 15–25\% \\
\bottomrule
\end{tabularx}
\sourcenote{本研究团队评估，折价率为相对于当前估值的风险调整幅度}
\end{exhibitbox}

经过风险调整后，估值排序发生了显著变化：时代电气的低估优势更加突出（风险最低），而天岳先进的估值吸引力大幅下降（40–50\% 的风险折价后，当前股价已不便宜）。这印证了一个核心原则——\textbf{高风险标的的"高增长预期"需要更大的安全边际来补偿}。

\begin{keyinsight}[估值审计核心结论]
经过系统性审计，我们对初始估值框架做以下修正：

1. \textbf{时代电气}估值最低且风险最低，经风险调整后仍是板块内最具性价比的标的，维持"核心配置"判断。
2. \textbf{扬杰科技}业绩质量最高，估值处于合理区间，经风险调整后仍具配置价值，维持"增持"。
3. \textbf{华润微、士兰微}当前 PE 看似很高，但属于周期底部效应，用 PB 和 EV/EBITDA 衡量估值更合理。经周期调整后估值并不极端，但需警惕复苏不及预期风险。
4. \textbf{天岳先进、三安光电}属于主题/困境投资范畴，PE 估值失效，PS/PB 估值缺乏可比锚点，经风险调整后安全边际不足，建议维持"高风险/关注"评级，仅适合风险承受能力强的投资者小仓位参与。
5. \textbf{斯达半导、新洁能}业绩兑现能力偏弱但估值不低，PEG 偏高，建议谨慎，等待业绩拐点或估值回调后的介入机会。
\end{keyinsight}
